\section{Conclusions}

This systematic literature review examined how fuzzy-logic-based inference systems are being used in real-time air-quality monitoring platforms for urban environments and associated risk assessment. The evidence shows an active and technically diverse field, with 19 primary studies that collectively demonstrate the viability of combining air-quality evaluation, fuzzy inference, and real-time operation. At the same time, the review confirms that the literature remains heterogeneous in its choice of input variables, fuzzy-system design, deployment architecture, and evaluation practices, which still limits direct comparability across studies.

Air-quality data and contextual variables show a clear pattern of concentration. PM$_{2.5}$ and CO are the most recurrent inputs, followed by CO$_2$ and PM$_{10}$, while temperature and humidity are the environmental variables most frequently used as contextual support. This pattern indicates that the reviewed platforms tend to privilege parsimonious and operationally viable input configurations, especially for AQI/IAQ classification, alerting, and pollutant prediction. At the same time, it suggests that the chemical and contextual coverage of the phenomenon often remains partial, particularly when systems are evaluated with only one or two pollutants or with limited reporting of supporting environmental variables.

Fuzzy-system design is dominated by configurations in which fuzzy inference occupies a central role in the decision process rather than acting as an auxiliary post-processing layer. Mamdani and ANFIS emerge as the dominant fuzzy-system configurations in the reviewed literature, although they are associated with different design profiles: Mamdani appears more often in interpretable classification and alerting settings, whereas ANFIS is more strongly associated with prediction and adaptive adjustment. Across the corpus, the main design dimensions that shape comparability are the purpose of inference, the definition of rules, the family of membership functions, the defuzzification strategy, the modeling of risk, and the explicit reporting of interpretability artifacts.

Real-time platform implementation is already a recurrent feature of the field, usually supported by IoT edge-cloud or two-layer architectures. Continuous monitoring, dashboard-based visualization, and, in some cases, alarm or actuation modules are already present in the reviewed literature. Nevertheless, end-to-end operational integration is still unevenly documented. In several studies, good model performance does not necessarily imply that acquisition, inference, publication, and response have been described with enough detail to support reproducibility or external transfer. The same applies to computational efficiency, calibration procedures, and long-term operational stability, which are reported much less systematically than predictive or classificatory performance.

The main research question can therefore be answered as follows: fuzzy-logic-based inference systems are being used primarily as central decision mechanisms in real-time air-quality monitoring platforms to classify AQI/IAQ states, estimate risk, support alerts, and, to a lesser extent, predict pollutant behavior. The reviewed studies provide consistent evidence of practical applicability, but the field remains fragmented across heterogeneous design choices and uneven reporting practices. As a result, the literature provides clear evidence of functional applicability, yet still only partial evidence of cross-context transferability, comparable evaluation, and fully reproducible operational deployment.

\subsection{Future Work}
Future work should address the main weaknesses identified in the review. First, new studies should strengthen multi-context validation, particularly across different cities, deployment conditions, and temporal windows, in order to improve external generalization. Second, future platforms would benefit from more explicit reporting of calibration procedures, hardware constraints, latency, and computational cost, since these dimensions are essential for reproducibility but remain underreported. Third, predictive systems should evaluate performance not only through aggregate error metrics but also through stability across forecast horizons. Finally, when fuzzy platforms are applied in healthcare or other risk-sensitive domains, future research should distinguish more clearly between operational utility, risk classification, and causal validation, so that preventive monitoring is not overstated as clinical evidence.\\
